% Seção 6 — Dados. Estado das fontes: CLAUDE.md (tabela de status) e
% docs/fontes-e-vintages.md. Sem descritivas de desfecho.

\chapter{Dados}
\label{cap:dados}

A unidade de análise é o município por mês de nascimento. O painel principal
cobrirá os municípios cearenses com aplicação aérea e os de comparação, de 2005
a 19 de dezembro de 2024; os 184 municípios do estado entram na forma reduzida,
e municípios de estados vizinhos, no controle sintético e na fronteira
(fontes no Apêndice~\ref{ap:fontes}).

\textbf{Desfechos.} Do SINASC serão obtidos o peso médio ao nascer, desfecho
confirmatório, e a proporção de nascidos com menos de 2.500 gramas e com menos
de 37 semanas de gestação, exploratórios (siglas no
Apêndice~\ref{ap:fontes}). A média é o desfecho confirmatório porque é nela
que o análogo mede o efeito e porque proporções de eventos raros são menos
precisas em células municipais pequenas; o baixo peso indica, sem teste, se o
efeito se concentra na cauda inferior. A série atravessa a troca do formulário de nascido vivo, em 2011
\cite{brasil2011sinasc}, que alterou o registro da idade gestacional; o peso ao
nascer não foi afetado do mesmo modo, e a prematuridade é analisada com essa
ressalva. Usa-se o município de residência da mãe. Do SIM será extraído o óbito
fetal, exploratório, e construído um desfecho composto, a proporção de
gestações registradas que terminam em nascido vivo com 2.500 gramas ou mais.
Do SINAN serão extraídas as notificações de intoxicação por agrotóxico
agrícola, por circunstância, para o canal de substituição
(Apêndice~\ref{ap:espec}), e as de arboviroses, para verificar se essa fonte
de sazonalidade do peso difere entre os grupos.

\textbf{Exposição.} A classificação pelo método vem dos registros da
Seção~\ref{sec:medicao} e, provisoriamente, do Censo Agropecuário de 2006; o
grupo de comparação, do Censo Agropecuário de 2017; a chuva, das séries de
precipitação por município. Como a lei estadual de 1993 autoriza os municípios
a legislar supletivamente, regras locais anteriores a 2019 alteram o período
de comparação; um levantamento preliminar encontrou uma proibição em Limoeiro
do Norte, de 20 de novembro de 2009 \cite{limoeiro2009lei1478}, revogada em 26
de maio de 2010 por lei que passou a regular a prática
\cite{limoeiro2010lei1511}. O levantamento será refeito, com conferência de
vigência.

\textbf{Escopo.} O canal hídrico, que exigiria medições de agrotóxicos na
água de consumo com localização dos pontos de coleta, não integra a análise
confirmatória, e a exposição não é medida no nível da residência. O §~2º do
art.~28-B alcança a dispersão aérea contra vetores; não se verificou se ela
ocorreu no Ceará antes de 2019, e, se ocorreu, há municípios tratados que a
classificação não vê.

\textbf{Reprodutibilidade e sigilo.} O processamento será feito por rotinas
versionadas, com data de extração registrada; microdados não são
redistribuídos, e células com poucos eventos não são divulgadas.

\textbf{Aspectos éticos.} Os dados de saúde são secundários, de acesso público
e sem identificação individual; pesquisas dessa natureza não são submetidas ao
sistema CEP/CONEP, nos termos da Resolução nº~510/2016 do Conselho Nacional de
Saúde \cite{cns2016res510}. Os registros obtidos pela Lei de Acesso à
Informação identificam empresas e responsáveis técnicos; serão usados apenas
para classificar municípios e meses, e nenhum resultado será divulgado com
identificação de empresa ou pessoa. Como os operadores são poucos, a medida
municipal e mensal pode identificá-los: agregação, limiar de supressão e
guarda dos documentos são fixados no plano de análise, antes do recebimento.
